% !TEX program = xelatex
\documentclass{CSUthesis}


\usepackage[figure,table]{totalcount}
\usepackage{totcount}

%参考文献计数
\newtotcounter{citnum}
\def\oldbibitem{} \let\oldbibitem=\bibitem
\def\bibitem{\stepcounter{citnum}\oldbibitem}
%\total{citnum}
%累积引用次数计数
%\newtotcounter{citesnum}
%\def\oldcite{} \let\oldcite=\cite
%\def\cite{\stepcounter{citesnum}\oldcite}
%\total{citesnum}

%%%%%%%%%%%%%%%%%%%%%%%%%%%%%%%%%%%%%%%%%%%%%%%%%%
% 前置部分的页眉页脚设置
% -----------------------------------------------%
% 正文和后置部分用阿拉伯数字编连续码，前置部分用罗马数字单独编连续码（封面除外）。
% 设置封面页后的页码

% 设置页眉和页脚 
\pagestyle{fancy}
% 正文以前部分无需页眉
\fancyhf{} % 清空原有格式
\renewcommand{\headrulewidth}{0pt}
% 中文摘要之前无需页码

%!TEX root = ../csuthesis_main.tex
% 文章信息
\titlecn{中南大学研究生学位论文LaTeX模板}
\titleen{LaTeX Template of Postgraduate Thesis of \\Central South University}

\priormajor{计算机科学与技术}
\minormajor{计算机应用技术}
\interestmajor{旁门左道}
\author{郭大侠}
\supervisor{我自己\ 教授}
\subsupervisor{}
\department{计算机学院}
\studentid{144601044}
\thesisdate{year=2020,month=4}



\clcnumber{TP391} 				% 中图分类号 Chinese Library Classification
\schoolcode{10533}			% 学校代码
\udc{004.9}						% UDC
\academiccategory{学术学位}	% 学术类别


\newif \ifblindreview % 条件语句，是否是盲审版本
\newif \ifAcademic % true为学术学位，false为专业学位
\newif \ifPublic   % true为公开，false为涉密



%\blindreviewtrue  %%盲审版本
\blindreviewfalse  %%正常版本

\Academictrue % 学术学位
%\Academicfalse % 专业学位
\Publictrue  % 公开




\begin{document}
%%%%%%%%%%%%%%%%%%%%%%%%%%%%%%%%%%%%%%%%%%%%%%%%%%
% 封面
% -----------------------------------------------%
\makecoverpage

\clearpage
\mbox{}
\clearpage

\ifblindreview	% 盲审不需要扉页和声明页
\else
%%%%%%%%%%%%%%%%%%%%%%%%%%%%%%%%%%%%%%%%%%%%%%%%%%
% 扉页 
% -----------------------------------------------%

\maketitlepage

\clearpage
\mbox{}
\clearpage

\setlength{\headheight}{9.6mm}
\setlength{\footskip}{7.9mm}
\setlength{\textheight}{681.5pt} % 调整版面高度
%%%%%%%%%%%%%%%%%%%%%%%%%%%%%%%%%%%%%%%%%%%%%%%%%%
% 声明页
% -----------------------------------------------%
\announcement
\clearpage
\mbox{}
\clearpage
\fi
%%%%%%%%%%%%%%%%%%%%%%%%%%%%%%%%%%%%%%%%%%%%%%%%%%

% 设置页眉和页脚 %
\setlength{\headheight}{9.6mm}
\setlength{\footskip}{7.9mm}
\pagestyle{fancy}
\fancyhf[CF]{\TimesNewRoman \zihao{-5} \thepage} % 所有（奇数和偶数）中间页脚，TimesNewRoman小五号

\addtocontents{toc}{\vspace{-8pt}}

% 中文摘要
% -----------------------------------------------%
\pagenumbering{Roman} % 重摘要至目录部分页码为大写罗马字体
\setcounter{page}{1} % 页码从I重新开始
\phantomsection                %% 解决目录中超链接地址错误问题
\addcontentsline{toc}{section}{摘要} %%增加摘要至目录且与第一章对齐
%!TEX root = ../csuthesis_main.tex
% 设置中文摘要
\keywordscn{中南大学；学位论文；LaTeX模板}
\categorycn{TP391}
\itemcountcn{图 \totalfigures\ 幅，表 \totaltables\ 个，参考文献 \total{citnum}\ 篇}

\begin{abstractcn}\setlength{\baselineskip}{20pt}%\renewcommand{\baselinestretch}{1.0}
LaTeX利用设置好的模板，可以编译为格式统一的pdf。目前国内大多出版社与高校仍在使用word，word由于其强大的功能与灵活性，在新手面对形式固定的论文时，排版、编号、参考文献等简单事务反而会带来很多困难与麻烦，对于一些需要通篇修改的问题，要想达到LaTeX的效率，对word使用者来说需要具有较高的技能水平。

为了能把主要精力放在论文撰写上，许多国际期刊和高校都支持LaTeX的撰写与提交，新手不需要关心格式问题，只需要按部就班的使用少数符号标签，即可得到符合要求的文档。且在需要全篇格式修改时，更换或修改模板文件，即可直接重新编译为新的样式文档，这对于word新手使用word的感受来说是不可思议的。

本项目的目的是为了创建一个符合中南大学研究生学位论文（博士）撰写规范的TeX模板，解决学位论文撰写时格式调整的痛点。
\end{abstractcn}
\clearpage

%%%%%%%%%%%%%%%%%%%%%%%%%%%%%%%%%%%%%%%%%%%%%%%%%%
% 英文摘要
% -----------------------------------------------%
\phantomsection                %% 解决目录中超链接地址错误问题
\addcontentsline{toc}{section}{ABSTRACT}
%!TEX root = ../csuthesis_main.tex
\keywordsen{CSU;~LaTeX;~Template}
\categoryen{TP391}
\begin{abstracten}
LaTeX can be compiled into a pdf of uniform format using the set template.At present, most domestic publishers and universities still use word. Because of its powerful function and flexibility, when faced with fixed-form papers by novices, simple matters such as typesetting, numbering, and reference documents will bring many difficulties and troubles. For some problems that need to be modified throughout, to achieve the efficiency of LaTeX, it requires a high level of skill for word users.

In order to focus on the writing of papers, many international journals and universities support the writing and submission of LaTeX. Novices don't need to care about formatting issues. They only need to use a few symbolic labels step by step to get the documents that meet the requirements. And when you need to modify the entire format, you can directly recompile the template file by replacing or modifying the template file. This is incredible for the word novice to use the word.

The purpose of this project is to create a TeX template that meets the specifications of the graduate degree thesis (PhD) of Central South University, and to address the pain points of format adjustment during the dissertation writing.
\end{abstracten}
\clearpage

% 目录
% -------------------------------------------%
{
	\renewcommand*{\baselinestretch}{1.3841}   % 行间距20pt
	\renewcommand{\contentsname}{\vspace{10pt} \hfill \heiti \bfseries \zihao{3} 目\quad 录\hfill ~\\}
	\phantomsection                %% 解决目录中超链接地址错误问题
	\addcontentsline{toc}{section}{目录} % 在目录中添加目录页码.
	\tableofcontents
}
\clearpage

% 插图索引
% -------------------------------------------%
{
	\setlength{\baselineskip}{20pt}         % 基准行间距
	\renewcommand{\baselinestretch}{1.0}   % 几倍行间距
	{~}
	\vspace{-10pt}
	\renewcommand{\listfigurename}{\hfill \heiti \zihao{3} 插图索引\hfill ~\\}  
	% 在目录中添加插图索引页码，按需确定是否注释下面两行
	\phantomsection                %% 解决目录中超链接地址错误问题
	\addcontentsline{toc}{section}{插图索引} 
	%%设置插图索引的 图 标签
	\let\oldnumberline\numberline%
	\renewcommand{\numberline}{\figurename~\oldnumberline}
	%增加图目录，按需确定是否注释下面一行
	\listoffigures
}

\clearpage

% 表格索引
% -------------------------------------------%
	{
	\setlength{\baselineskip}{20pt}         % 基准行间距
	\renewcommand{\baselinestretch}{1.0}   % 几倍行间距
	{~}
	\vspace{-10pt}
	\renewcommand{\listtablename}{\hfill \heiti \zihao{3} 表格索引\hfill ~\\}  
	 % 在目录中添加表格索引页码，按需确定是否注释下面两行
	 \phantomsection                %% 解决目录中超链接地址错误问题
	\addcontentsline{toc}{section}{表格索引}
	%%设置表格索引的 表 标签
	\let\oldnumberline\numberline%
	\renewcommand{\numberline}{\tablename~\oldnumberline}
	%增加表目录，按需确定是否注释下面一行
	\listoftables
}

\clearpage


%%%%%%%%%%%%%%%%%%%%%%%%%%%%%%%%%%%%%%%%%%%%%%%%%%
% 符号说明（必要时）
% -----------------------------------------------%
\phantomsection                %% 解决目录中超链接地址错误问题
\addcontentsline{toc}{section}{符号说明} % 在目录中添加目录页码.
\begin{denotation}
\vspace{6pt} \songti \zihao{5}
%\songti \zihao{5}
\begin{center}%此处最好是h
\arrayrulewidth1pt
\resizebox{\textwidth}{!}{
\begin{tabular}{| m{0.4\textwidth}<{\centering} | m{0.3\textwidth}<{\centering} | m{0.3\textwidth}<{\centering} |}
\hline
符号&意义&单位（量纲）\\
\hline
频率&赫[兹]&Hz\\
\hline
\end{tabular}}
\end{center}
\end{denotation}


%%%%%%%%%%%%%%%%%%%%%%%%%%%%%%%%%%%%%%%%%%%%%%%%%%
% 正文页眉页脚
% -----------------------------------------------%

\renewcommand{\headrulewidth}{1pt}

% 去掉页眉章节序号后面的“.”
\renewcommand{\sectionmark}[1]{\markboth{第{\thesection}章~ #1}{第{\thesection}章~ #1}}
\renewcommand{\subsectionmark}[1]{\markright{\leftmark}}
\renewcommand{\subsubsectionmark}[1]{\markright{\leftmark}}

\fancyhf[RH]{\songti \zihao{5} \leftmark} % 设置所有（奇数和偶数）右侧页眉net
\ifDoctor
	\fancyhf[LH]{\songti \zihao{5} 中南大学博士学位论文} % 设置所有（奇数和偶数）左侧页眉
\else
	\fancyhf[LH]{\songti \zihao{5} 中南大学硕士学位论文} % 设置所有（奇数和偶数）左侧页眉
\fi
% 正文内容 
% --------------------------------------------%
\setcounter{page}{1} % 重置页码编号
\pagenumbering{arabic} % 设置页码编号为阿拉伯数字

% 可以使用include命令导入tex文件，从而避免过多修改本文件。

% 论文正文是主体，主体部分应从另页右页开始，每一章应另起页。一般由序号标题、文字叙述、图、表格和公式等五个部分构成。

% 重新设置正文行间距，因为前置部分设置时候行间距被改过
\renewcommand*{\baselinestretch}{1.0}   % 几倍行间距
\setlength{\baselineskip}{20pt}         % 基准行间距

% 正文
{
	% 表格字号应比正文小，一般五号/10.5pt，但是暂时没法再cls里设置（不然会影响到封面等tabular环境）
	% 所以目前只好在主文件里局部\AtBeginEnvironment
	\AtBeginEnvironment{tabular}{\zihao{5}}
	%!TEX root = ../csuthesis_main.tex
% 论文正文是主体，主体部分应从另页右页开始，每一章应另起页。一般由序号标题、文字叙述、图、表格和公式等五个部分构成。
\section{绪论}
\subsection{研究背景与意义}

目的是创建一个符合中南大学研究生学位论文（博士）撰写规范的LaTeX模板，解决学位论文撰写时格式调整的痛点。

已有珠玉在前，我们之所以还要重新造轮子，主要是参考了2020年4月2号学校下发的[《中南大学研究生学位论文撰写规范》中大研字【2020】30号]( http://oa.its.csu.edu.cn/Home/ReleaseMainText/909CD53BF50943CD97B83C352032FEA4 )，重新修改了页面布局、字体类型和大小、标题内容，以期做到与 Word 模板尽可能的相似。主要修改如下：
\begin{itemize}
\item 依据专家意见修订定理和证明等环境，如”定理”使用小四黑体，编号随章节变化重新编号（如定理4-1），定理内容使用小四宋体，且内容行距与正文一致；”证明”无需编号，且以黑色小方块结尾；
\item 修订算法在每个章节重新编号问题；
\item 增加符号说明页和附录页（如果不需要，请在.cls文件对应处注释掉即可）；
\item  增加参考文献按国标 gbt7714-2015要求，只核对了常用的图书、中英文期刊，会议格式，其余未常使用的未进行核对（如有问题请改回gbt7714-2005）；
\item  修订多个子图Caption居中问题；
\item 依据专家意见调整成果与致谢部分间距，并增加目录中的点密度；
\item 按照图书馆最新要求（2020年12月份），去除目录中红色边框；
\item 增加页眉信息：中南大学博士论文与右侧的章节名保持一致，以及无需号章节名保持一致;
\item 增加中英文摘要至目录，并保持与章节名昨对其;
\item 参考文献完全依照国标 gbt7714-2005，修正了部分 Bug，提供了新的引用命令；
\item 按照最新版本要求，在声明扉页前后各增加一页空白页，保证装订单独成页；
\item 章节标题居中，并改成‘第1章’样式；
\item 目录中，将原章节标题换成‘第几章’样式，字体按要求加粗；
\item 中文摘要到目录结束用罗马数字编写页码，小五号Times New Roman,居中；
\item 增加插图索引和表格索引；
\item 所有的章节题目和中英文摘要均按要求修改字体和间距；
\end{itemize}

\subsection{主要研究工作}
我堂堂双一流高校竟然没有官方研究生论文LaTeX模板！！！虽然我LaTeX水平也很水……但是通过大量debug也勉强给大家凑出来一个格式绝对标准的LaTeX模板，模板代码丑就丑吧，能用就行。写了大量注释，有一点LaTeX基础就可以根据自己需要修改CSUthesis.cls文件。

(1) 提供图片插入示例。

(2) 提供表格插入示例。

(3) 提供公式插入示例。

(4) 提供参考文献插入示例。

\subsection{论文组织结构}

全文内容共六章，具体内容组织如下：

第一章为绪论。

第二章为图片插入示例。

第三章为表格插入示例。

第四章为公式插入示例。

第五章为参考文献插入示例。

第六章总结与展望，总结了本文的主要工作，展望了下一阶段的研究方向。

\newpage

\section{图像布局}


\subsection{单图布局}



\textbf{单图布局如图\ref{F.csu_single}所示。}

\begin{figure}[hbt]
\centering
\includegraphics[width=0.5\textwidth]{csu.png}
\caption{单图布局示例}
\label{F.csu_single}
\end{figure}

\subsection{横排布局}

\textbf{横排布局如图\ref{F.csu_row}所示。}

\begin{figure}[!htb]
    \centering
    \begin{subfigure}[t]{0.24\linewidth}
    	\captionsetup{justification=centering} %子图caption居中
        \begin{minipage}[b]{1\linewidth}
        \includegraphics[width=1\linewidth]{csu.png}
         \caption{}
        \end{minipage}
    \end{subfigure}
    \begin{subfigure}[t]{0.24\linewidth}
    \captionsetup{justification=centering} %子图caption居中
        \begin{minipage}[b]{1\linewidth}
        \includegraphics[width=1\linewidth]{csu.png}
        \caption{}
        \end{minipage}
    \end{subfigure}
    \begin{subfigure}[t]{0.24\linewidth}
    	\captionsetup{justification=centering} %子图caption居中
        \begin{minipage}[b]{1\linewidth}
        \includegraphics[width=1\linewidth]{csu.png}
        \caption{}
        \end{minipage}
    \end{subfigure}
    \begin{subfigure}[t]{0.24\linewidth}
    	\captionsetup{justification=centering} %子图caption居中
        \begin{minipage}[b]{1\linewidth}
        \includegraphics[width=1\linewidth]{csu.png}
        \caption{}
        \end{minipage}
    \end{subfigure}
    \caption{横排布局示例}
    \label{F.csu_row}
\end{figure}



\subsection{竖排布局}
\textbf{竖排布局如图\ref{F.csu_col}所示。}

\begin{figure}[!htb]
    \centering
    \begin{subfigure}[t]{0.15\linewidth}
        \begin{minipage}[b]{1\linewidth}
        \includegraphics[width=1\linewidth]{csu.png}
        \caption{}
        \end{minipage}
    \end{subfigure}\\
    \begin{subfigure}[t]{0.15\linewidth}
        \begin{minipage}[b]{1\linewidth}
        \includegraphics[width=1\linewidth]{csu.png}
        \caption{}
        \end{minipage}
    \end{subfigure}
    \caption{竖排布局示例}
    \label{F.csu_col}
\end{figure}



\subsubsection{竖排多图横排布局}

\begin{figure}[!htb]
    \centering
    \begin{subfigure}[t]{0.13\linewidth}
    	\captionsetup{justification=centering} %子图caption居中
        \begin{minipage}[b]{1\linewidth}
        \includegraphics[width=1\linewidth]{csu.png} \vspace{-1ex} \vfill
        \includegraphics[width=1\linewidth]{csu.png}
         \caption{}
        \end{minipage}
    \end{subfigure}
    \begin{subfigure}[t]{0.13\linewidth}
    	\captionsetup{justification=centering} %子图caption居中
        \begin{minipage}[b]{1\linewidth}
        \includegraphics[width=1\linewidth]{csu.png} \vspace{-1ex} \vfill
        \includegraphics[width=1\linewidth]{csu.png}
        \caption{}
        \end{minipage}
    \end{subfigure}
    \caption{竖排多图横排布局}
    \label{F.csu_col_row}
\end{figure}

\textbf{竖排多图横排布局如图\ref{F.csu_col_row}所示。注意看(a)、(b)编号与图关系。}


\subsubsection{横排多图竖排布局}



\begin{figure}[!htb]
    \centering
    \begin{subfigure}[t]{0.3\linewidth}
    	\captionsetup{justification=centering} %子图caption居中
        \begin{minipage}[b]{1\linewidth}
        \includegraphics[width=0.45\linewidth]{csu.png}
        \includegraphics[width=0.45\linewidth]{csu.png}
        \caption{}
        \end{minipage}
    \end{subfigure}\\
    \begin{subfigure}[t]{0.3\linewidth}
    	\captionsetup{justification=centering} %子图caption居中
        \begin{minipage}[b]{1\linewidth}
        \includegraphics[width=0.45\linewidth]{csu.png}
        \includegraphics[width=0.45\linewidth]{csu.png}
        \caption{}
        \end{minipage}
    \end{subfigure}
    \caption{横排多图竖排布局，斜体emph \emph{A}，A，斜体texit \textit{A}}
    \label{F.csu_row_col}
\end{figure}

\textbf{横排多图竖排布局如图\ref{F.csu_row_col}所示。注意看(a)、(b)编号与图关系。}

\subsection{本章小结}
本章示例图片布局。

\newpage


\section{表格插入示例}

\begin{table}[htb]
  \centering
  \caption{表格为三线表斜体emph \emph{A}，A，斜体texit \textit{A}}
  \label{T.example}
  \begin{tabular}{llllll}
  \toprule
    & \emph{A}A \textit{A}  & B  & C  & D  & E \\
  \midrule
1 	& 212 & 414 & 4 		& 23 & fgw	\\
2 	& 212 & 414 & v 		& 23 & fgw	\\
3 	& 212 & 414 & vfwe		& 23 & 嗯	\\
4 	& 212 & 414 & 4fwe		& 23 & 嗯	\\
5 	& af2 & 4vx & 4 		& 23 & fgw	\\
6 	& af2 & 4vx & 4 		& 23 & fgw	\\
7 	& 212 & 414 & 4 		& 23 & fgw	\\
\bottomrule

\end{tabular}
\end{table}

\textbf{表格如表\ref{T.example}所示，latex表格技巧很多，这里不再详细介绍。}



\newpage

\section{算法示例}


\begin{algorithm}  
	\caption{Fourier-Mellin Based KCF}  
	\label{alg:A}  
	\hspace*{0.02in}{\bf Input:}
	Image $I$\\preprocessed kernelized template $T_\kappa$\\
	\hspace*{0.02in}{\bf Output:} 
	scale $\sigma$, angle $\theta$ relation between $I$ and $T$ 
	
	\begin{algorithmic}[1] 
		\STATE{fourier transform: $F=\mathcal{F}(I)$}
		\STATE {high pass filter: $F_h=\mathcal{H}(F)$\\$\mathcal{H}(x,y)=(1.0-cos(\pi x)cos(\pi y))(2.0-cos(\pi x)cos(\pi y))$}
		\STATE {log-polar transform: $F_{lp}=\mathcal{L}(F_h)$}
		\STATE {apply kernel function: $F_\kappa=\mathcal{K}(F_{lp})$}
		\STATE {phase correlation: $(\Delta x, \Delta y)=\mathcal{C}(F_\kappa, T_\kappa)$}
		\STATE {resolove scale and rotation:\\
			$\theta=\alpha \Delta x$, $\sigma=log(\Delta y)$\\
			where $\alpha$ is translation factor of pixel translation on fourier domain and polar angle on origin image
		}
	\end{algorithmic}  
\end{algorithm}


  \begin{algorithm}
	\caption{算法示例}
	\label{alg:SPSH01}
	\renewcommand{\algorithmicrequire}{\textbf{Input:}}
	\renewcommand{\algorithmicensure}{\textbf{Output:}}
	
	\begin{algorithmic}[1]
		\REQUIRE 相关输入。。。。
		
		\ENSURE 相关输出。。。
		
		\STATE 算法描述  % 只占用一个行号
		\FOR{$i \gets 1\cdots N$}
		\STATE 算法描述
		\FOR{\textbf{each} $j \gets 1\cdots K$}
		\STATE 算法描述
		\ENDFOR
		\ENDFOR
		\REPEAT
		\REPEAT
		\STATE 令$\tau\gets\tau+1$
		\UNTIL 内循环迭代终止条件
		\STATE 。。。
		\UNTIL 外循环迭代终止条件
	\end{algorithmic}
\end{algorithm}

\textbf{如算法\ref{alg:A}所示，latex算法技巧很多。按需调整，这里不再详细介绍。}

\newpage

\section{公式、定理、证明插入示例}

\begin{flalign}
\text{P1: } &\min_{\eta,R_u>0,R_d>0}\big\{ T_{\text{latency}}(\eta,R_u,R_d)\big\}  \\
&\text{s.t. }   0 \leq \eta \leq 1 \label{P1C1} 
\end{flalign}

\textbf{公式插入示例如公式(\ref{E.example})所示。}

\begin{equation}
\gamma_{x}=
\left\{
  \begin{array}{lr}
  0, & {\rm if}~~\;|x| \leq \delta \\
  x, & {\rm otherwise}
  \end{array}
\right.
\label{E.example}
\end{equation}

\begin{flalign}
&\text{P1:} \max_{\bigg\{\substack{P_{m,i},P_{n,i}\\
		q_{m,i},q_{n,i}\\ \forall n,m,i}\bigg\}}\big[R_{\text{sum}}(P_{m,i},P_{n,i},q_{m,i},q_{n,i},\forall n,m,i)\big],  \\
& \text{s.t.}~~q_{m,i}\in(0,1), q_{n,i}\in(0,1), \forall n,m,i, \label{allocons1} \\
& \hspace{1.5em} ~0\leq \sum_{i=1}^Rq_{n,i}P_{n,i}\leq P_n^{sum}, \forall n, \label{powercons1}  \\ 
& \hspace{1.5em} ~0\leq \sum_{i=1}^Rq_{m,i}P_{m,i}\leq P_m^{sum}, \forall m, \label{powercons2} \\ 
& \hspace{1.5em} ~\sum_{i=1}^Rq_{n,i}\leq 1,\sum_{i=1}^Rq_{m,i}\leq 1,\forall m,n, \label{RBcons} \\ 
& \hspace{1.5em} ~\sum_{n=1}^Nq_{n,i}\leq 1,\sum_{m=1}^Mq_{m,i}\leq 1,\forall i, \label{RBcons2} \\
& \hspace{1.5em} ~C_{m,BS,i}(P_{m,i},q_{m,i})\geq\varepsilon_{m,i},\forall m, \label{capcons}
\end{flalign}


  %%%%%%%%%%%%%% 示例 %%%%%%%%%%%%%%%%
公式子编号示例：
\begin{subequations}\label{Eq.(3-4)}
	\renewcommand{\theequation}{\theparentequation-\alph{equation}} % 定义编号方式
	\begin{align}
		&\varphi_{n,t}\in\left\{0,1\right\}, \forall n\in\mathcal{N}, t\in\mathcal{N}_{\mathrm{t}}, \label{Eq.(3-4a)} \\
		&\varphi_{n,t}\in\left\{0,1\right\}, \forall n\in\mathcal{N}, t\in\mathcal{N}_{\mathrm{t}}, \label{Eq.(3-4b)} \\
		&\varphi_{n,t}\in\left\{0,1\right\}, \forall n\in\mathcal{N}, t\in\mathcal{N}_{\mathrm{t}}, \label{Eq.(3-4c)}
	\end{align}
\end{subequations}

其中，公式\ref{Eq.(3-4a)}表示。

  %%%%%%%%%% 示例 %%%%%%%%%%%5
定理和证明环境说明：如”定理”使用小四黑体，编号随章节变化重新编号（如定理4-1），定理内容使用小四宋体，且内容行距与正文一致；”证明”无需编号，且以黑色小方块结尾。

\begin{theorem}\label{theorem1}
	\setlength{\baselineskip}{20pt}         % 基准行间距
	\renewcommand{\baselinestretch}{1.0}   % 几倍行间距
	开始定理。。。
\end{theorem}

\begin{proof*} %% 证明不编号
	\setlength{\baselineskip}{20pt}         % 基准行间距
	\renewcommand{\baselinestretch}{1.0}   % 几倍行间距
	开始证明。。。
	$\hfill\blacksquare$ %% 以黑方块结尾
\end{proof*}

\newpage

\section{参考文献插入示例}

LaTeX\cite{lamport1994latex}插入参考文献最方便的方式是使用bibliography\cite{pritchard1969statistical}，大多数出版商的论文页面都会有导出bib格式参考文献的链接，建议使用Jabref管理参考文献，把每个文献的bib放入``thesis-references''，然后用bibkey即可插入参考文献。

中文文献\cite{zh-book-1}，注意手动编辑bibkey为英文的即可。

可以将文献标注为右上角\citess{shiweisong2019},只需要在现有的cite后加“ss”即可。

英文会议\cite{Kraus2021Current}, \cite{WuYangLuEtAl2021}.

英文期刊\cite{LuoZengYuanEtAl2016}, \cite{Wu2022Boosting}.


\textbf{特别强调：}从Google下载的bib也不一定全是对的，如发现有信息缺失，请下载原文核对。比如已发表的期刊，要包保证年、卷、标。

\textbf{注意：}如发现替换后的参考文献没有更新，请删除主文件夹下xxx.bbl文件，重新编译即可。

\newpage


\section{总结与展望}

\noindent{纯数字编号}
\begin{enumerate}
 \item XXXXXXXXXX
 \label{item1}
 \item XXXXXXXXXX
 \item XXXXXXXXXX
\end{enumerate}
罗马编号
\begin{enumerate}[label=(\roman*)]
 \item XXXXXXXXXX
 \label{item2}
 \item XXXXXXXXXX
 \item XXXXXXXXXX
\end{enumerate}
括号编号
\begin{enumerate}[label=(\arabic*)]
 \item XXXXXXXXXX
 \label{item3}
 \item XXXXXXXXXX
 \item XXXXXXXXXX
\end{enumerate}
半括号编号
\begin{enumerate}[label=\arabic*)]
 \item XXXXXXXXXX
 \label{item4}
 \item XXXXXXXXXX
 \item XXXXXXXXXX
\end{enumerate}
小字母编号
\begin{enumerate}[label=\alph*)]
 \item XXXXXXXXXX
 \label{item5}
 \item XXXXXXXXXX
 \item XXXXXXXXXX
\end{enumerate}

引用测试,正如\ref{item1}、\ref{item2}、\ref{item3}、\ref{item4}、\ref{item5}所示

\subsection{工作展望}
手动编号 %（不推荐,无法被交叉引用）
\par
本课题针对XX，鉴于XXX，对XX进行了提高，但是XXX，所以有如下XX：

（1）目前XX虽然XX，但是XX仍然XX，所以XX仍然是一个值得XX的问题。

（2）随着XX，XX具有XX的问题，仍值得进一步XX。

（3）本课题在XX有了XX，但是XX的XX还存在XX，所以XX。


\newpage

}

	%%%%%%%%%%%%%%%%%%%%%%%%%%%%%%%%%%%%%%%%%%%%%%%%%%
	% 临时标签，用于编译时追踪正文末尾
	%%%%%%%%%%%%%%%%%%%%%%%%%%%%%%%%%%%%%%%%%%%%%%%%%%
	
	%%%%%%%%%%%%%%%%%%%%%%%%%%%%%%%%%%%%%%%%%%%%%%%%%%
	% 后续内容，标题三号黑体居中，章节无编号
	% --------------------------------------------%
	
	% https://www.zhihu.com/question/29413517/answer/44358389 %
	% 说明如下：
	% secnumdepth 这个计数器是 LaTeX 标准文档类用来控制章节编号深度的。在 article 中，这个计数器的值默认是 3，对应的章节命令是 \subsubsection。也就是说，默认情况下，article 将会对 \subsubsection 及其之上的所有章节标题进行编号，也就是 \part, \section, \subsection, \subsubsection。LaTeX 标准文档类中，最大的标题是 \part。它在 book 和 report 类中的层级是「-1」，在 article 类中的层级是「0」。这里，我们在调用 \appendix 的时候将计数器设置为 -2，因此所有的章节命令都不会编号了。不过，一般还是会保留 \part 的编号的。所以在实际使用中，将它设置为 0 就可以了。

% 在修改过程中请注意不要破环命令的完整性

\renewcommand\appendix{\setcounter{secnumdepth}{-2}}
\appendix

% 主文件有代码去掉页眉章节编号的“.”，但这会因为bug导致无编号章节显示一个错误编号，所以这里在无编号章节之前再次重定义sectionmark。
\renewcommand{\sectionmark}[1]{\markboth{#1}{#1}}
\renewcommand{\subsectionmark}[1]{\markright{\leftmark}}
\renewcommand{\subsubsectionmark}[1]{\markright{\leftmark}}

% section 标题从这里往后改为三号黑体居中
\titleformat{\section}{\centering \zihao{3} \bfseries \heiti}{\thesection}{1em}{}
\titlespacing*{\section} {0pt}{18pt}{32pt}

% \section{参考文献} % bibliography会自动显示参考文献四个字

% \nocite{*} % 该命令用于显示全部参考文献，即使文中没引用
% cls文件中已经引入package，这里不需要调用 \bibliographystyle 了。
% \bibliographystyle{gbt7714-2005} 
{~}
\vspace{18pt}
\phantomsection                                      %% 解决目录中超链接地址错误问题
\addcontentsline{toc}{section}{参考文献} % 由于参考文献不是section，这句把参考文献加入目录
\zihao{-4} \setmainfont{Times New Roman}
\setlength{\baselineskip}{20pt}          % 基准行间距 
\bibliographystyle{gbt7714/2015-numerical}  % 使用2005会带来一些错误
\bibliography{references}

\clearpage
% 附录

\titleformat{\section}{\centering \zihao{3} \bfseries \heiti}{\thesection}{1em}{}
\titlespacing*{\section} {0pt}{18pt}{32pt}
\titleformat{\subsection}{\zihao{-4} \songti}{\thesubsection}{1em}{}
\titlespacing*{\subsection} {2em}{0pt}{0pt}
\titleformat{\subsubsection}{\zihao{-4} \songti}{\thesubsubsection}{1em}{}
\titlespacing*{\subsubsection} {2em}{0pt}{0pt}

{~}
\vspace{18pt}
\section{附录A （附录名称 ）（三号黑体，加粗）（必要时）} % 无章节编号

附录正文…
…（格式参考正文）。
换行示例
。

%\newpage
%
%{~}
%\vspace{-9pt}
%\section{附录B （附录名称 ）（三号黑体，加粗）（必要时）} % 无章节编号
%
%附录正文…
%…（格式参考正文）。
%换行示例
%。


\clearpage
% 研究成果和致谢

\titleformat{\section}{\centering \zihao{3} \bfseries \heiti}{\thesection}{1em}{}
\titlespacing*{\section} {0pt}{18pt}{32pt}
\titleformat{\subsection}{\zihao{-4} \songti}{\thesubsection}{1em}{}
\titlespacing*{\subsection} {2em}{0pt}{0pt}
\titleformat{\subsubsection}{\zihao{-4} \songti}{\thesubsubsection}{1em}{}
\titlespacing*{\subsubsection} {2em}{0pt}{0pt}

%!TEX root = ../csuthesis_main.tex

{~}
\vspace{18pt}
\section{攻读学位期间主要的研究成果} % 无章节编号

\ifblindreview
% \noindent
% （盲审隐去作者相关具体信息）
\fi
\vspace{11pt}
\subsection*{一、学术论文}

\ifblindreview

% \noindent
% 第一作者：JCR 1 区 x 篇，会议 x 篇 \\{}
% 第二作者：JCR 1 区 x 篇，3 区 x 篇，4 区 x 篇，EI x 篇 

% \noindent
% 投稿状态： 
% IEEE Transactions on Image Processing 1篇（under review）\\{}
% IEEE Transactions on Circuits and Systems for Video Technology 1 篇（Accept with Minor Revision）

% 学位办老师要求用如下这种几乎算是单盲的格式，我也木有办法……
\subsubsection*{已录用/检索论文}
%\noindent
第一作者：
\begin{enumerate}[label={[\arabic*]},itemindent=2em,wide]
	\item CSU Latex Template[J]. CSU player: 1(1):1-10. {\bfseries \heiti（SCI 录用，JCR 1 区）}
	\item CSU Latex Template[J]. CSU player: 1(1):1-10. {\bfseries \heiti（SCI 检索，JCR 2 区）}
\end{enumerate}
第二作者：
\begin{enumerate}[label={[\arabic*]},itemindent=2em,wide]
	\item CSU Latex Template[J]. CSU player: 1(1):1-10. {\bfseries \heiti（SCI 检索，JCR 1 区）}
	\item CSU Latex Template[J]. CSU player: 1(1):1-10. {\bfseries \heiti（SCI 检索，JCR 2 区）}
\end{enumerate}
%\noindent
\subsubsection*{投稿状态论文}
%\noindent
第一作者：
\begin{enumerate}[label={[\arabic*]},itemindent=2em,wide]
	\item CSU Latex Template. XXX Transactions on CSU player. {\bfseries \heiti（SCI Under Review，JCR 1 区）}
\end{enumerate}
第二作者：
\begin{enumerate}[label={[\arabic*]},itemindent=2em,wide]
	\item CSU Latex Template. XXX Transactions on CSU player. {\bfseries \heiti（SCI Under Review，JCR 1 区）}
	\item CSU Latex Template. XXX Transactions on CSU player. {\bfseries \heiti（SCI Under Review，JCR 2 区）}
\end{enumerate}


\else
% 标准版本
\begin{enumerate}[label={[\arabic*]},itemindent=2em,wide]
	\item \textbf{Daxia Mou}, Director, Someone. CSU Latex Template[J]. CSU player: 1(1):1-10. {\bfseries \heiti（SCI 检索，JCR 1 区）}
	\item Director, \textbf{Daxia Mou}, Someone, Someother. XXXXXX[J]. Transactions on Image Processing. {\bfseries \heiti（SCI Under Review，JCR 1 区）}
	\item Director, \textbf{Daxia Mou}, Someone, Someother. XXXXXX[J]. Transactions on Circuits and Systems for Video Technology. {\bfseries \heiti（SCI Under Review，JCR 1 区）}
\end{enumerate}
\fi

\vspace{22pt}
\subsection*{二、发明专利}
\ifblindreview
发明专利 2 项，已公开
\else
\begin{enumerate}[label={[\arabic*]},itemindent=2em,wide]
	\item 某大侠，XXX，XXX. 一种用Latex写中南大学学位论文的方法. 申请号：CN20190415xxxx，公开号：CNXXXXXXXXXA
\end{enumerate}
\fi

\ifblindreview
\else

\vspace{22pt}
\subsection*{三、主持和参与的科研项目}
\begin{enumerate}[label={[\arabic*]},itemindent=2em,wide]
	\item 国家自然科学基金面上项目《XXXXXXXXXXXX》， 项目编号：XXXXXXXX，参与.
\end{enumerate}

% {~}
\vspace{22pt}
\subsection*{四、个人获奖情况}
%\noindent
\begin{enumerate}[label={[\arabic*]},itemindent=2em,wide]
	\item XX金奖
	\item XX奖学金
\end{enumerate}
\fi

\newpage

\ifblindreview
\else

{~}
\vspace{18pt}
\section{{致~~~~谢}} % 无章节编号
% 重新设置正文行间距，因为前置部分设置时候行间距被改过
\renewcommand*{\baselinestretch}{1.0} % 几倍行间距
\setlength{\baselineskip}{20pt} % 基准行间距

作者对给予指导、各类资助和协完成研究工以及提供种论文有 作者对给予指导、各类资助和协完成研究工以及提供种论文有 利条件的单位及个人表示感谢。

致谢应实事求是，切忌浮夸与庸俗之词 。

\newpage
\fi



\end{document}
